% !TeX root = ../教材题库_全部条目_转写工作稿.tex
% Source: XJ-B1-C04-HW-P125-Y1398
% 第4章 / 4.3 对数函数 / 教材 p.125、126 / PDF 第132、133页
% OCR未校对：公式与文字均可能有误；未完成数学TeX转写；禁止直接发给学生
% 原题及相邻原解截图（核对用）：../题目原图/XJ-B1-C04-HW-P125-Y1398-01.webp, ../题目原图/XJ-B1-C04-HW-P125-Y1398-02.webp
% 本文件只在转写完成后作为题干使用；OCR 条目尚不是可靠题干。
\begin{bankitem}{习题4.3 第12题}{第4章 / 4.3 对数函数 / 教材 p.125、126 / PDF 第132、133页}{XJ-B1-C04-HW-P125-Y1398}{OCR未校对：公式与文字均可能有误；未完成数学TeX转写；禁止直接发给学生}
\begin{ocrdraft}
[教材p.125 / PDF第132页]
\par
12．设山乃，0都是不等于1的正数，函数
\par
y==logax，3，==log/,x，3'一，y==logdX．在同
\par
Y-zogeX
\par
直角坐标系中的图象如右图所示，则a，乃，。d的
\par
y=log\ocrglyph{201E}x
\par
大小关系是
\par
（第12题）
\par
[教材p.126 / PDF第133页]
\par
温故而知新
\end{ocrdraft}
\end{bankitem}
